\chapter{The Multi-Agent Path Finding Problem}
\label{ch:ch07}
% Function names for pseudocode are prefixed with "Mapf" to stay unique
% across the book.
\SetKwFunction{MapfPosition}{Position}
\SetKwFunction{MapfPairwise}{FirstConflictPairwise}
\SetKwFunction{MapfHashed}{FirstConflictHashed}

\begin{objectives}
\begin{itemize}
  \item State the classical multi-agent path finding problem exactly: graph, agents, discrete synchronous time, move and wait actions, time-indexed paths and plans, and the stay-at-target convention.
  \item Recognise the five kinds of conflict of Stern et al.\ (vertex, edge, swapping, following, cycle), know which two this book forbids and why, and detect them by hand in a table of paths.
  \item Compute the sum of costs and the makespan of a plan, explain why the two objectives can disagree, and state what is known about the complexity of feasibility and of optimality.
  \item Read and write instances in the \code{.map} and \code{.scen} formats of the standard benchmark and name the solver families that the following chapters develop.
  \item Implement a plan validator and a conflict detector in \python, in $\bigO{k^2T}$ and in $\bigO{kT}$ time, and use it to build deliberate conflict cases.
\end{itemize}
\end{objectives}

% ---------------------------------------------------------------------
\section{Why this matters for a drone swarm}
\label{sec:ch07-motivation}

Twelve inspection drones share the hall of a warehouse. Each has its own start
pad and its own target shelf, and each, on its own, can find a shortest route
with the \astar of \cref{ch:ch04} in a few milliseconds. Fly those twelve
routes at the same time and two of them meet head-on in an aisle, a third
crosses the first at an intersection at the wrong moment, and a fourth arrives
at a shelf in front of which a drone that has already finished is still
hovering. None of the single-agent planners of
\cref{ch:ch03,ch:ch04,ch:ch05,ch:ch06} can see these events, because each of
them plans one agent in a world that contains no others.

\Cref{fig:ch07-conflict-growth} measures how bad the problem is. On random
$16 \times 16$ grids with $20\,\%$ obstacles we planned every agent
independently and counted the conflicts of the resulting \textbf{naive
plan}\index{naive plan}. With two agents, $95\,\%$ of the plans happen to be
conflict-free. With eight, only $5\,\%$ are, and a plan contains $3.5$
conflicts on average; with twenty agents no naive plan out of sixty was free
of conflicts, and with forty agents a plan has about $97$ conflicts among
$72$ conflicting pairs. The number of conflicts grows roughly like the number
of pairs, $k(k-1)/2$. Coordination is not an optional refinement: beyond a
handful of agents, independent planning never produces a flyable plan.

\begin{figure}[tb]
  \centering
  \inputfigure{ch07/conflict-growth}
  \caption{Naive plans (independent shortest paths) on random $16 \times 16$
  grids with $20\,\%$ obstacles, $60$ instances per value of $k$
  (\code{gen\_ch07\_conflict\_growth.py}). Left: the mean number of vertex and
  swapping conflicts and of conflicting agent pairs grows roughly
  quadratically with $k$. Right: the fraction of instances whose naive plan is
  conflict-free drops to zero at about twelve agents.}
  \label{fig:ch07-conflict-growth}
\end{figure}

The problem of finding routes for many agents that never collide is
\textbf{multi-agent path finding}\index{MAPF (multi-agent path finding)}
(\textbf{MAPF}). This chapter defines it; the four chapters that follow solve
it. In the hybrid architecture of \cref{ch:ch24}, MAPF is the job of the
\emph{global} layer: it produces the nominal, time-indexed routes of all
drones under our control before take-off, and it is consulted again whenever
the local avoidance layer or the replanning layer has changed a route and the
swarm must re-check that the changed route conflicts with nobody. That
re-check is the conflict detector of this chapter, and it is the one piece of
MAPF that every layer of the system uses.

A problem chapter has a different shape from an algorithm chapter.
\Cref{sec:ch07-intuition} gives the idea and \cref{sec:ch07-problem} the exact
definitions of Stern et al.~\cite{stern2019mapf}: instances, actions, paths,
plans, the five kinds of conflict, and the two objectives. The algorithm of
the chapter, in \cref{sec:ch07-detection}, is conflict detection, and
\cref{sec:ch07-example} runs it by hand on three agents.
\Cref{sec:ch07-properties} collects what is known about complexity:
feasibility is easy, optimality is NP-hard for both objectives.
\Cref{sec:ch07-solvers} maps the solver families to
\cref{ch:ch08,ch:ch09,ch:ch10,ch:ch11,ch:ch22}, \cref{sec:ch07-benchmarks}
explains the benchmark files you will use for the rest of Part~III,
\cref{sec:ch07-implementation} the \python, and \cref{sec:ch07-drones} what
changes when the agents are drones.

% ---------------------------------------------------------------------
\section{The idea in plain words}
\label{sec:ch07-intuition}

Freeze time into steps of length $\dt$ and write down where every agent is at
every step. A route becomes a row of a table, one column per time step, and a
plan for $k$ agents is a table with $k$ rows. Two agents \emph{conflict} when
the table shows them in the same place in the same column, or shows them
exchanging places along one edge between two consecutive columns. Finding a
plan means filling in the table so that every row is a legal route for its
agent and no such coincidence occurs anywhere; finding a good plan means doing
so with rows that are as short as possible.

This is the space-time view of \cref{ch:ch04}: the state of an agent is a pair
$(v, t)$, a step either moves along an edge or waits, and time advances by one
for every agent at once. What is new is that the rows are no longer
independent. The same table also shows what a conflict is \emph{not}. If
agent~2 leaves a cell in the step in which agent~1 enters it, the two never
share a column entry; they follow each other, one cell apart, like the
carriages of a train. Whether that is dangerous depends on the physics of the
agents (\cref{sec:ch07-drones}), but under the discrete model it is not a
conflict. \Cref{fig:ch07-conflict-types} shows the five patterns that the
literature distinguishes.

\begin{figure}[tb]
  \centering
  \inputfigure{ch07/conflict-types}
  \caption{The conflict types of Stern et al.\ \cite{stern2019mapf} between
  time $t$ and $t+1$. (a)~Two agents enter the same vertex. (b)~Two agents
  traverse the same edge in the same direction, which means they already
  shared $u$ at time $t$. (c)~Two agents exchange vertices along an edge.
  (d)~One agent enters the vertex that the other is leaving. (e)~A rotation
  along a fully occupied cycle. Classical MAPF forbids (a)--(c) and allows (d)
  and (e).}
  \label{fig:ch07-conflict-types}
\end{figure}

\begin{keyidea}
A MAPF plan is a table of positions, one row per agent and one column per
time step, in which an agent that has reached its goal keeps its goal cell in
every later column. The plan is valid when every row is a legal path and no
two rows share a cell in a column or swap two cells between consecutive
columns. Everything else, including how good a plan is and how to find one,
is built on this table.
\end{keyidea}

% ---------------------------------------------------------------------
\section{Problem statement and notation}
\label{sec:ch07-problem}

\subsection{Instances, actions and paths}

\begin{definition}[Classical MAPF instance]
\label{def:ch07-instance}
A \textbf{classical MAPF instance}\index{MAPF!classical}\index{MAPF!instance}
is a triple $\langle G, s, g \rangle$: an undirected graph $G = (V, E)$, a
number $k$ of \textbf{agents}\index{agent} $a_1, \dots, a_k$, and two maps
$s, g \colon \set{1,\dots,k} \to V$ that assign to every agent a
\textbf{start vertex} $s_i$ and a \textbf{goal vertex} $g_i$. The starts are
pairwise distinct and so are the goals. Time is discrete,
$t \in \set{0, 1, 2, \dots}$, and all agents act synchronously: in every time
step each agent either \textbf{moves}\index{move action} along one edge to an
adjacent vertex or \textbf{waits}\index{wait action} at its vertex. Both
actions take exactly one time step and cost $1$.
\end{definition}

The graph is usually the 4-connected grid $G_4$ of \cref{ch:ch02}, whose
vertices are the free cells and whose edges join horizontally or vertically
adjacent free cells. We write a cell as $(r, c)$ with the row $r$ counted from
the top and the column $c$ from the left, because that is the order in which
a benchmark map file is read (\cref{sec:ch07-benchmarks}); the single-agent
chapters used $(x, y)$ with the origin at the bottom left, and both are only
names for vertices. A start vertex may coincide with a goal vertex, of the
same agent or of another one; only the starts among themselves and the goals
among themselves must be distinct.

\begin{definition}[Time-indexed path, stay-at-target]
\label{def:ch07-path}
A \textbf{time-indexed path}\index{path!time-indexed} for agent $a_i$ is a
sequence $\pi_i = (\pi_i[0], \pi_i[1], \dots, \pi_i[T_i])$ of vertices with
$\pi_i[0] = s_i$, $\pi_i[T_i] = g_i$ and, for every $t < T_i$, either
$\pi_i[t+1] = \pi_i[t]$ (a wait) or $\set{\pi_i[t], \pi_i[t+1]} \in E$ (a
move). Under the \textbf{stay-at-target}\index{stay-at-target} convention the
agent remains at its goal after the end of its path: we set
$\pi_i[t] = g_i$ for all $t > T_i$, so that $\pi_i[t]$ is defined for every
$t \ge 0$. The \textbf{cost}\index{path!cost (MAPF)} of the path is the time
of its last arrival at the goal,
\begin{equation}
  \cost(\pi_i) = \min \set{ \tau \ge 0 \;:\; \pi_i[t] = g_i \text{ for all } t \ge \tau }.
  \label{eq:ch07-path-cost}
\end{equation}
\end{definition}

Two remarks. First, a path may visit its goal, leave it and come back;
\cref{eq:ch07-path-cost} charges every step until the \emph{final} arrival,
and trailing waits at the goal are free. The sequence
$\bigl((0,0), (0,1), (0,1), (0,2), (0,2)\bigr)$ has $T_i = 4$ but cost $3$.
Second, the alternative \textbf{disappear-at-target}\index{disappear-at-target}
convention removes an agent from the graph when it first reaches its goal.
It models robots that park in a docking station or leave the warehouse, and
it makes some instances easier because a finished agent blocks nobody. A
hovering drone blocks its cell as much as a flying one, so this book uses
stay-at-target throughout, as do the benchmark and almost all solvers
\cite{stern2019mapf}.

\begin{definition}[Plan and solution]
\label{def:ch07-plan}
A \textbf{plan}\index{plan!MAPF} is a tuple $\Pi = (\pi_1, \dots, \pi_k)$ of
one time-indexed path per agent. Its \textbf{horizon}\index{plan!horizon} is
$T = \max_i T_i$. A plan is a \textbf{solution}\index{MAPF!solution} (or is
\emph{valid}) if it contains no vertex conflict and no swapping conflict in
the sense of \cref{def:ch07-conflicts} at any time $t \ge 0$.
\end{definition}

\subsection{Conflicts}
\label{sec:ch07-conflicts}

\begin{definition}[Conflicts \cite{stern2019mapf}]
\label{def:ch07-conflicts}
Let $\Pi$ be a plan, $i \ne j$ two agents and $t \ge 0$ a time step.
\begin{enumerate}[label=(\alph*)]
  \item A \textbf{vertex conflict}\index{conflict!vertex} $\langle a_i, a_j, v, t \rangle$
    occurs if $\pi_i[t] = \pi_j[t] = v$: both agents occupy $v$ at time $t$.
  \item An \textbf{edge conflict}\index{conflict!edge} $\langle a_i, a_j, u, v, t \rangle$
    occurs if $\pi_i[t] = \pi_j[t] = u$ and $\pi_i[t+1] = \pi_j[t+1] = v \ne u$:
    both traverse the edge $\set{u,v}$ in the same direction in the same step.
  \item A \textbf{swapping conflict}\index{conflict!swapping} $\langle a_i, a_j, u, v, t \rangle$
    occurs if $\pi_i[t] = \pi_j[t+1] = u$ and $\pi_i[t+1] = \pi_j[t] = v \ne u$:
    the agents exchange vertices along the edge $\set{u,v}$ between $t$ and $t+1$.
  \item A \textbf{following conflict}\index{conflict!following} occurs if
    $\pi_i[t+1] = \pi_j[t]$: agent $a_i$ enters at $t+1$ the vertex that $a_j$
    occupied at $t$.
  \item A \textbf{cycle conflict}\index{conflict!cycle} occurs if agents
    $a_{i_1}, \dots, a_{i_m}$ ($m \ge 2$) all move at step $t$ and
    $\pi_{i_1}[t+1] = \pi_{i_2}[t]$, $\pi_{i_2}[t+1] = \pi_{i_3}[t]$, \dots,
    $\pi_{i_m}[t+1] = \pi_{i_1}[t]$: they rotate along a cycle of $m$ vertices
    that is fully occupied.
\end{enumerate}
\end{definition}

The five types are not independent, and it pays to see exactly how they
overlap before deciding which of them to forbid.

\begin{proposition}[Relations between conflict types]
\label{prop:ch07-hierarchy}
\begin{enumerate}[label=(\alph*)]
  \item An edge conflict $\langle a_i, a_j, u, v, t\rangle$ implies the vertex
    conflicts $\langle a_i, a_j, u, t\rangle$ and $\langle a_i, a_j, v, t+1\rangle$.
  \item A swapping conflict is a cycle conflict with $m = 2$, and it consists of
    two following conflicts, one in each direction.
  \item A cycle conflict with $m \ge 3$ consists of $m$ following conflicts and
    contains no vertex conflict and no swapping conflict.
  \item A following conflict in which the leading agent waits,
    $\pi_j[t+1] = \pi_j[t] = \pi_i[t+1]$, is the vertex conflict
    $\langle a_i, a_j, \pi_j[t], t+1\rangle$.
\end{enumerate}
\end{proposition}
\begin{proof}
Parts (a), (b) and (d) are the definitions read twice. For (c), the $m$
vertices of the cycle are distinct and at time $t$ each holds exactly one of
the agents; at $t+1$ each agent has moved to the vertex of its predecessor, so
the $m$ agents again occupy the $m$ distinct vertices and no vertex conflict
arises. A swap would require $\pi_{i_1}[t+1] = \pi_{i_2}[t]$ together with
$\pi_{i_2}[t+1] = \pi_{i_1}[t]$, that is $i_3 = i_1$, which is impossible for
$m \ge 3$.
\end{proof}

\paragraph{What this book forbids.}
Following \textcite{stern2019mapf}, \emph{classical} MAPF forbids vertex and
swapping conflicts and allows following and cycle conflicts. Vertex conflicts
are collisions by definition. Swapping conflicts are collisions too: two
agents that exchange the endpoints of an edge in one step pass through each
other in the middle of it. Edge conflicts need no rule of their own, because
by \cref{prop:ch07-hierarchy}(a) they are already vertex conflicts. Following
is allowed because synchronous unit-speed motion keeps the follower one full
edge behind the leader at every instant, and forbidding it would rule out
convoys through a corridor and the plan of \cref{sec:ch07-example}. Cycle
conflicts are allowed for the same reason: a rotation on a fully occupied
cycle keeps every agent one edge from its neighbours. Some hardware disagrees
with both choices, for example robots that need a free cell ahead before they
start moving, or drones so large that following around a corner brings them
too close (\cref{exr:ch07-separation}); such systems add following conflicts
to the forbidden list and pay with longer plans.
\Cref{tab:ch07-conflict-summary} sums up the rules.

\begin{table}[tb]
  \centering\small
  \caption{The conflict types of \cref{def:ch07-conflicts} and their status in classical MAPF.}
  \label{tab:ch07-conflict-summary}
  \begin{tabular}{@{}llll@{}}
  \toprule
  Conflict & Pattern between $t$ and $t+1$ & Classical MAPF & Why\\
  \midrule
  vertex & $\pi_i[t] = \pi_j[t]$ & forbidden & two agents in one cell\\
  edge & same edge, same direction & forbidden & implies a vertex conflict\\
  swapping & same edge, opposite directions & forbidden & agents pass through each other\\
  following & $\pi_i[t+1] = \pi_j[t]$, $a_j$ moves on & allowed & one edge apart at all times\\
  cycle & rotation on an occupied cycle & allowed & $m$ following conflicts, no collision\\
  \bottomrule
  \end{tabular}
\end{table}

\subsection{Objectives}
\label{sec:ch07-objectives}

\begin{definition}[Sum of costs and makespan]
\label{def:ch07-objectives}
The \textbf{sum of costs}\index{sum of costs}\index{objective!sum of costs}
and the \textbf{makespan}\index{makespan}\index{objective!makespan} of a plan
$\Pi = (\pi_1, \dots, \pi_k)$ are
\begin{equation}
  \sumcost(\Pi) = \sum_{i=1}^{k} \cost(\pi_i),
  \qquad
  \makespan(\Pi) = \max_{i=1,\dots,k} \cost(\pi_i).
  \label{eq:ch07-objectives}
\end{equation}
A solution is \textbf{optimal}\index{MAPF!optimal solution} for an objective
if no solution of the instance has a smaller value of it.
\end{definition}

The sum of costs is the total number of time steps that agents spend before
they are finished, so it measures total flight time and energy; the makespan
is the time at which the last agent arrives, so it measures the duration of
the mission. Every wait before the final arrival counts in both. The two
objectives agree on what a good plan looks like but not on which plan is
best, and it takes only three agents to separate them.

\begin{example}[Where the objectives disagree]
\label{ex:ch07-disagree}
\Cref{fig:ch07-objectives-disagree} shows three agents on an empty
$4 \times 5$ grid: $a_1$ flies along row~1 from $(1,0)$ to $(1,4)$, $a_2$ down
column~1 from $(0,1)$ to $(2,1)$, and $a_3$ up column~2 from $(3,2)$ to
$(0,2)$. Their shortest paths have costs $4$, $2$ and $3$ and collide twice:
$a_1$ and $a_2$ both reach $(1,1)$ at $t=1$, and $a_1$ and $a_3$ both reach
$(1,2)$ at $t=2$. Plan~$Y$ lets $a_1$ wait one step at its start; it then
reaches each crossing one step after the other agent has passed, the costs
are $5, 2, 3$, and $\sumcost(Y) = 10$, $\makespan(Y) = 5$. Plan~$X$ instead
lets $a_2$ and $a_3$ each wait once, so that $a_1$ passes first: costs
$4, 3, 4$, $\sumcost(X) = 11$, $\makespan(X) = 4$. A search of the joint
state space (\code{optimal\_soc\_bruteforce} and
\code{optimal\_makespan\_bruteforce} in \code{ch07\_mapf.py}) confirms that
$10$ and $4$ are the optimal values and that no solution attains both: every
solution with $\sumcost = 10$ has makespan $5$.
\end{example}

\begin{figure}[tb]
  \centering
  \inputfigure{ch07/objectives-disagree}
  \caption{Three agents whose shortest paths cross at $(1,1)$ and $(1,2)$.
  Plan $Y$ delays the long path once and minimises the sum of costs; plan $X$
  delays the two short paths and minimises the makespan. No plan achieves
  $\sumcost = 10$ and $\makespan = 4$ at the same time.}
  \label{fig:ch07-objectives-disagree}
\end{figure}

Which objective should a drone swarm use? The sum of costs is the standard
choice of the MAPF literature and of \cbs (\cref{ch:ch09}), because it
spreads delays over the agents and because it decomposes over agents, which
the solvers exploit. The makespan is the natural choice when the mission ends
only when the last drone has landed. \Cref{ch:ch25} reports both.

\subsection{Feasibility versus optimality}

An instance is \textbf{feasible}\index{MAPF!feasibility} (or solvable) if it
has at least one solution. Feasibility can fail for simple reasons, such as a
goal that is unreachable from its start, and for combinatorial ones: two
agents at the ends of a corridor that must exchange places have no solution
unless the corridor has a pocket to step into. Optimality asks for the best
solution with respect to $\sumcost$ or $\makespan$. \Cref{sec:ch07-properties}
shows that the two questions are of very different difficulty, and
\cref{sec:ch07-solvers} sorts the solvers by which of them they answer and
how well.

% ---------------------------------------------------------------------
\section{The algorithm: conflict detection}
\label{sec:ch07-detection}

The solvers of \cref{ch:ch08,ch:ch09,ch:ch10,ch:ch11} all contain the same
subroutine: given a plan, find its first conflict or report that there is
none.\index{conflict!detection} It is also the validator that you run on any
plan before you trust it, and the re-check of the drone architecture.
\Cref{alg:ch07-pairwise} states the direct version.

\begin{algorithm}[htb]
\caption{First conflict of a plan by pairwise comparison}
\label{alg:ch07-pairwise}
\KwIn{plan $\Pi = (\pi_1, \dots, \pi_k)$, each $\pi_i$ a list of $T_i + 1$ vertices}
\KwOut{the earliest vertex or swapping conflict, or \emph{nil}}
\Fn{\MapfPosition{$\pi_i$, $t$}}{
  \lIf{$t \le T_i$}{\KwRet $\pi_i[t]$}
  \KwRet $\pi_i[T_i]$\tcp*{stay-at-target padding}
}
\Fn{\MapfPairwise{$\Pi$}}{
  $T \gets \max_i T_i$\label{alg:ch07-pairwise:horizon}\;
  \For{$t \gets 0$ \KwTo $T$}{
    \ForEach{pair $i < j$}{
      \If{\MapfPosition{$\pi_i$, $t$} $=$ \MapfPosition{$\pi_j$, $t$}}{\KwRet vertex conflict $\langle a_i, a_j, \pi_i[t], t\rangle$\label{alg:ch07-pairwise:vertex}\;}
    }
    \If{$t < T$}{
      \ForEach{pair $i < j$}{
        $u \gets$ \MapfPosition{$\pi_i$, $t$}; $v \gets$ \MapfPosition{$\pi_i$, $t+1$}\;
        \If{$u \ne v$ \KwAnd \MapfPosition{$\pi_j$, $t$} $= v$ \KwAnd \MapfPosition{$\pi_j$, $t+1$} $= u$}{\KwRet swapping conflict $\langle a_i, a_j, u, v, t\rangle$\label{alg:ch07-pairwise:swap}\;}
      }
    }
  }
  \KwRet \emph{nil}\;
}
\end{algorithm}

\paragraph{Walkthrough.}
The helper \MapfPosition implements the padding of \cref{def:ch07-path}: a
question about a time beyond the end of a path is answered with the goal.
This is what makes an agent that has arrived block its cell for everybody who
comes later, and it is the first thing that home-made detectors get wrong
(\cref{sec:ch07-implementation}). Line~\ref{alg:ch07-pairwise:horizon} fixes
the horizon. Scanning further is pointless: after time $T$ no agent moves, so
no swap can occur, and two agents that share a cell at some $t > T$ already
shared it at time $T$; since goals are distinct, two parked agents never
share a cell at all. For each time step the algorithm first looks for vertex
conflicts (line~\ref{alg:ch07-pairwise:vertex}) and then, unless $t = T$, for
swaps between $t$ and $t+1$ (line~\ref{alg:ch07-pairwise:swap}). The
condition $u \ne v$ excludes a pair of waiting agents, which is a vertex
conflict and has been reported already. Returning at the first hit gives the
earliest conflict, with ties broken towards vertex conflicts and then towards
the smallest pair $(i, j)$; dropping the returns and collecting the hits
instead gives all conflicts in the same order.

\paragraph{Cost.}
Each time step compares all $k(k-1)/2$ pairs twice, so the running time is
$\bigO{k^2 T}$. For a solver that calls the detector at every node of a search
tree this quadratic factor matters, and it is unnecessary: at a fixed time,
the agents at a vertex $v$ are found by hashing positions, not by comparing
pairs. \Cref{alg:ch07-hashed} keeps, per time step, a dictionary from vertices
to the agent seen there, and a second dictionary from directed edges
$(u, v)$ to the agent traversing them; a swap is a traversal of $(v, u)$ by
an earlier agent when $(u, v)$ is inserted.

\begin{algorithm}[htb]
\caption{First conflict of a plan with hash maps}
\label{alg:ch07-hashed}
\KwIn{plan $\Pi$ as in \cref{alg:ch07-pairwise}}
\KwOut{the earliest vertex or swapping conflict, or \emph{nil}}
\Fn{\MapfHashed{$\Pi$}}{
  $T \gets \max_i T_i$\;
  \For{$t \gets 0$ \KwTo $T$}{
    $\mathit{occupied} \gets$ empty map from vertices to agents\;
    \For{$i \gets 1$ \KwTo $k$}{
      $v \gets$ \MapfPosition{$\pi_i$, $t$}\;
      \lIf{$v \in \mathit{occupied}$}{\KwRet vertex conflict $\langle a_{\mathit{occupied}[v]}, a_i, v, t\rangle$}
      $\mathit{occupied}[v] \gets i$\label{alg:ch07-hashed:occ}\;
    }
    \lIf{$t = T$}{\KwRet \emph{nil}}
    $\mathit{moving} \gets$ empty map from directed edges to agents\;
    \For{$i \gets 1$ \KwTo $k$}{
      $u \gets$ \MapfPosition{$\pi_i$, $t$}; $v \gets$ \MapfPosition{$\pi_i$, $t+1$}\;
      \If{$u \ne v$}{
        \lIf{$(v, u) \in \mathit{moving}$}{\KwRet swapping conflict $\langle a_{\mathit{moving}[(v,u)]}, a_i, v, u, t\rangle$}
        $\mathit{moving}[(u, v)] \gets i$\label{alg:ch07-hashed:mov}\;
      }
    }
  }
  \KwRet \emph{nil}\;
}
\end{algorithm}

\begin{proposition}[Correctness and cost of conflict detection]
\label{prop:ch07-detection}
\Cref{alg:ch07-pairwise,alg:ch07-hashed} return \emph{nil} if and only if the
plan is a solution in the sense of \cref{def:ch07-plan}, and otherwise a
conflict with the smallest possible time $t$. \Cref{alg:ch07-pairwise} runs
in $\bigO{k^2 T}$ time; \cref{alg:ch07-hashed} runs in $\bigO{kT}$ expected
time and uses $\bigO{k}$ extra space.
\end{proposition}
\begin{proof}
Both algorithms test, for each $t \le T$, exactly the conditions of
\cref{def:ch07-conflicts}(a) and (c) on the padded paths, and the walkthrough
showed that no conflict of either kind can first occur after $T$. In
\cref{alg:ch07-hashed} a vertex conflict at $t$ between $a_j$ and $a_i$ with
$j < i$ is detected when $a_i$ is inserted, because $a_j$ was stored at
line~\ref{alg:ch07-hashed:occ} earlier in the same step; a swap between $a_j$
and $a_i$ is detected at the insertion of the later of the two directed
edges. Each time step performs $\bigO{k}$ dictionary operations of expected
constant time, and the maps are cleared every step.
\end{proof}

When several conflicts share a time step, the two versions may report
different ones; the implementation of \cref{sec:ch07-implementation} breaks
ties by the smallest agent pair in both, so that they agree, and its self-test
checks the agreement on random plans.

% ---------------------------------------------------------------------
\section{A worked example}
\label{sec:ch07-example}

\begin{example}[Three agents on a $4 \times 4$ grid]
\label{ex:ch07-grid}
\Cref{fig:ch07-example-instance} shows a $4 \times 4$ grid whose cells
$(1,0)$ and $(1,2)$ are blocked. Agent $a_1$ must fly along the top row from
$(0,0)$ to $(0,3)$, agent $a_2$ along the same row in the opposite direction
from $(0,3)$ to $(0,0)$, and agent $a_3$ starts in the pocket $(1,1)$ below
the row and wants the cell $(0,1)$ right above it. The instance is
\code{worked\_example()} in \code{ch07\_mapf.py}; every number below is
checked by its self-test, which numbers the agents from $0$ where the text
numbers them from $1$.
\end{example}

\begin{figure}[tb]
  \centering
  \inputfigure{ch07/example-instance}
  \caption{\Cref{ex:ch07-grid} with the naive plan of independent shortest
  paths. Solid circles are starts, dashed rings goals, and the small numbers
  give the time at which each agent reaches a cell. The crosses mark the three
  conflicts found by hand in \cref{tab:ch07-example-paths}: two vertex
  conflicts at $(0,1)$ and a swap on the edge $(0,1)$--$(0,2)$.}
  \label{fig:ch07-example-instance}
\end{figure}

\paragraph{The naive plan.}
Planning each agent alone with breadth-first search (or with \astar,
\cref{ch:ch04}) gives the three paths in the upper half of
\cref{tab:ch07-example-paths}: $\pi_1$ and $\pi_2$ run along the top row in
three moves each and $\pi_3$ needs a single move up. Their costs are $3$,
$3$ and $1$, so $\sumcost = 7$ and $\makespan = 3$; by \cref{thm:ch07-bounds}
below no solution can beat either number. Agent $a_3$ arrives at $t = 1$, and
the grey entries are the padding of \cref{def:ch07-path}: from then on
$\pi_3[t] = (0,1)$.

\begin{table}[tb]
  \centering\small
  \caption{The naive plan (top) and the conflict-free plan (bottom) of
  \cref{ex:ch07-grid} as tables of positions. Grey entries are positions
  after the final arrival, added by the stay-at-target padding; they cost
  nothing, but they do block the cell.}
  \label{tab:ch07-example-paths}
  \begin{tabular}{@{}lccccccr@{}}
  \toprule
  Agent & $t=0$ & $t=1$ & $t=2$ & $t=3$ & $t=4$ & $t=5$ & $\cost(\pi_i)$\\
  \midrule
  \multicolumn{8}{@{}l}{\emph{Naive plan (independent shortest paths): $\sumcost = 7$, $\makespan = 3$, three conflicts}}\\
  $a_1$ & $(0,0)$ & $(0,1)$ & $(0,2)$ & $(0,3)$ & \textcolor{sbGray}{$(0,3)$} & \textcolor{sbGray}{$(0,3)$} & 3\\
  $a_2$ & $(0,3)$ & $(0,2)$ & $(0,1)$ & $(0,0)$ & \textcolor{sbGray}{$(0,0)$} & \textcolor{sbGray}{$(0,0)$} & 3\\
  $a_3$ & $(1,1)$ & $(0,1)$ & \textcolor{sbGray}{$(0,1)$} & \textcolor{sbGray}{$(0,1)$} & \textcolor{sbGray}{$(0,1)$} & \textcolor{sbGray}{$(0,1)$} & 1\\
  \midrule
  \multicolumn{8}{@{}l}{\emph{Conflict-free plan, optimal for both objectives: $\sumcost = 12$, $\makespan = 5$}}\\
  $a_1$ & $(0,0)$ & $(0,1)$ & $(1,1)$ & $(0,1)$ & $(0,2)$ & $(0,3)$ & 5\\
  $a_2$ & $(0,3)$ & $(0,2)$ & $(0,1)$ & $(0,0)$ & \textcolor{sbGray}{$(0,0)$} & \textcolor{sbGray}{$(0,0)$} & 3\\
  $a_3$ & $(1,1)$ & $(2,1)$ & $(2,1)$ & $(1,1)$ & $(0,1)$ & \textcolor{sbGray}{$(0,1)$} & 4\\
  \bottomrule
  \end{tabular}
\end{table}

\paragraph{Detection by hand.}
Run \cref{alg:ch07-pairwise} on the upper table. The horizon is $T = 3$. At
$t = 0$ the three positions $(0,0)$, $(0,3)$, $(1,1)$ are distinct, and no
pair reverses an edge between $t = 0$ and $t = 1$. At $t = 1$ the column
reads $(0,1)$, $(0,2)$, $(0,1)$: agents $a_1$ and $a_3$ share $(0,1)$, and
the algorithm returns the vertex conflict $\langle a_1, a_3, (0,1), 1\rangle$.
If we collect every conflict instead of stopping, the swap check between
$t = 1$ and $t = 2$ finds $a_1$ moving $(0,1) \to (0,2)$ while $a_2$ moves
$(0,2) \to (0,1)$, the swapping conflict
$\langle a_1, a_2, (0,1), (0,2), 1\rangle$. At $t = 2$ the column reads
$(0,2)$, $(0,1)$, $(0,1)$, a second vertex conflict
$\langle a_2, a_3, (0,1), 2\rangle$, caused by nothing but the padding:
$a_3$ has been parked at its goal since $t = 1$ and $a_2$ runs into it. The
moves between $t = 2$ and $t = 3$ reverse no edge, the column at $t = 3$
holds three different cells, and the scan ends. The plan has exactly these
three conflicts, in this order, and
\code{all\_conflicts(naive\_plan(inst))} prints the same three.

\paragraph{The space-time view.}
\Cref{fig:ch07-spacetime} draws the same three paths in the time-expanded
graph of \cref{ch:ch04}\index{space-time!MAPF view}: one copy of every vertex
per time step, a vertical edge for a wait and a diagonal edge for a move. A
vertex conflict is a space-time node used by two paths; a swapping conflict
is a pair of crossing move edges. This picture is behind every solver of
\cref{ch:ch08,ch:ch09,ch:ch10}. A single agent is planned by space-time
\astar; a vertex constraint $\langle a_i, v, t\rangle$ deletes the node
$(v, t)$ for agent $a_i$, an edge constraint $\langle a_i, u, v, t\rangle$
deletes its move edge from $(u, t)$ to $(v, t+1)$, and the goal-stay test
makes sure that the agent can park at $g_i$ for ever once it arrives. What a
MAPF solver adds is the choice of the constraints; what this chapter adds is
the detector that tells it which conflict to resolve next.

\begin{figure}[tb]
  \centering
  \inputfigure{ch07/spacetime}
  \caption{The naive plan of \cref{ex:ch07-grid} in space-time. Columns are
  the cells of the top row and the pocket $(1,1)$, rows are time steps, and
  the grey lines are the wait and move edges of the time-expanded graph. The
  two red rings on nodes are the vertex conflicts at $(0,1)$; the ring between
  the rows marks the crossing move edges of the swap.}
  \label{fig:ch07-spacetime}
\end{figure}

\paragraph{A conflict-free plan.}
The lower half of \cref{tab:ch07-example-paths} and
\cref{fig:ch07-example-solution} show a solution. Agent $a_3$ backs down into
$(2,1)$ and waits there for one step; $a_1$ steps into the vacated pocket
$(1,1)$ at $t = 2$ and lets $a_2$ pass along the row; then $a_1$ climbs back
and continues, and $a_3$ follows one step behind it to its goal. Check the
columns: at $t = 2$ the positions are $(1,1)$, $(0,1)$, $(2,1)$, at $t = 3$
they are $(0,1)$, $(0,0)$, $(1,1)$, no column repeats a cell, and no pair of
moves reverses an edge. The plan does contain following conflicts. Between
$t = 1$ and $t = 2$, $a_2$ enters $(0,1)$ in the step in which $a_1$ leaves
it; between $t = 2$ and $t = 3$, $a_1$ enters $(0,1)$ as $a_2$ leaves it while
$a_3$ enters $(1,1)$ as $a_1$ leaves it, a convoy of three. Classical MAPF
allows all of them, and \code{validate\_plan} accepts the plan. Its costs are
$5$, $3$ and $4$, so $\sumcost = 12$ and $\makespan = 5$: coordination costs
five extra time steps in total and two extra steps of mission time compared
with the lower bounds $7$ and $3$. The reason is easy to see. Both $a_1$ and
$a_2$ need the whole top row and one of them must step aside; the only side
cells are $(1,1)$ and $(1,3)$, and $(1,3)$ is useless because whoever hides
there has to come back through $(0,3)$, where $a_1$ parks. So the pocket
$(1,1)$ must be used, and $a_3$ has to clear it first. The brute-force
searches of the joint state space in \code{ch07\_mapf.py} confirm that no
solution has $\sumcost < 12$ or $\makespan < 5$, so this plan is optimal for
both objectives at once; \cref{ex:ch07-disagree} showed that this is a
property of the instance, not a rule.

\begin{figure}[tb]
  \centering
  \inputfigure{ch07/example-solution}
  \caption{The conflict-free plan of \cref{ex:ch07-grid} as six snapshots;
  an arrow shows the move that an agent executes in the step that follows.
  At $t = 2$, $a_1$ hides in the pocket while $a_2$ passes; from $t = 3$ on,
  $a_3$ follows $a_1$ one cell behind, which classical MAPF allows.}
  \label{fig:ch07-example-solution}
\end{figure}

% ---------------------------------------------------------------------
\section{Properties: feasibility, optimality and complexity}
\label{sec:ch07-properties}

\subsection{Lower bounds}

\begin{proposition}[Elementary bounds]
\label{thm:ch07-bounds}
Let $d(u, v)$ be the length of a shortest path between $u$ and $v$ in $G$
and let $\Pi$ be any solution of $\langle G, s, g\rangle$. Then
\begin{equation}
  \sumcost(\Pi) \ge \sum_{i=1}^{k} d(s_i, g_i), \qquad
  \makespan(\Pi) \ge \max_{i} d(s_i, g_i), \qquad
  \makespan(\Pi) \le \sumcost(\Pi) \le k \cdot \makespan(\Pi).
  \label{eq:ch07-bounds}
\end{equation}
\end{proposition}
\begin{proof}
The path $\pi_i$ is at $g_i$ from time $\cost(\pi_i)$ on and makes at most
one move per step, so $\cost(\pi_i) \ge d(s_i, g_i)$; summing or maximising
over $i$ gives the first two inequalities. The third holds because a sum of
$k$ non-negative numbers lies between their maximum and $k$ times their
maximum.
\end{proof}

The first two bounds are the values of the naive plan\index{naive plan!lower
bound}, and they are the starting point of every optimal solver: the root of
the constraint tree of \cbs (\cref{ch:ch09}) is the naive plan, the bound is
what \ecbs (\cref{ch:ch10}) compares its solutions against, and the gap
between the bound and the optimal cost is the price of coordination. In
\cref{ex:ch07-grid} that gap is $5$ for the sum of costs and $2$ for the
makespan.

\subsection{Feasibility is easy}

\begin{theorem}[Feasibility can be decided in polynomial time]
\label{thm:ch07-feasible}
Let $G$ be connected and let at least two vertices be empty,
$k \le \abs{V} - 2$. Whether $\langle G, s, g\rangle$ has a solution can be
decided in polynomial time, and if it has one, a solution with
$\bigO{\abs{V}^3}$ moves can be constructed in polynomial
time~\cite{kornhauser1984pebble,dewilde2014push}.
\end{theorem}
\begin{proof}[Proof sketch]
Kornhauser, Miller and Spirakis~\cite{kornhauser1984pebble} studied the
\textbf{pebble motion}\index{pebble motion} problem, the sequential version
of MAPF in which one agent at a time moves to an adjacent empty vertex. They
showed that its solvability can be decided in polynomial time and that a
solvable instance has a solution of $\bigO{\abs{V}^3}$ moves. A sequence of
single moves is a classical MAPF plan in which all other agents wait, so a
solvable pebble instance is a solvable MAPF instance. Conversely, as long as
one vertex is empty, every parallel step of a MAPF solution can be
serialised. The moves of a step form chains and cycles; a chain ends at an
empty vertex and is executed from its head backwards. For a rotation on a
fully occupied cycle, walk the empty vertex to a neighbour of the cycle by a
sequence of single moves, step one agent of the cycle out into it, move the
remaining agents of the cycle one by one, step the agent back in at the
other end, and walk the empty vertex back along the same route, which undoes
the displacement of the agents on that route. The two problems therefore
have the same solvable instances, and the decision procedure for pebble
motion decides MAPF. Push-and-Rotate (\cref{ch:ch11}) turns this argument
into an algorithm: it is complete for instances with at least two empty
vertices, reports the unsolvable ones, and builds a solution in polynomial
time~\cite{dewilde2014push}.
\end{proof}

Two remarks put the theorem in perspective. First, solvable does not mean
easy to see. Two agents at the ends of a corridor that must exchange places
have no solution however many empty cells the corridor contains, because on
a path graph the left-to-right order of the agents can only change through a
swap (\cref{exr:ch07-corridor}); a single pocket cell makes the instance
solvable. With exactly one empty vertex, the classic case is the
15-puzzle\index{15-puzzle}, where half of all arrangements of the tiles are
unreachable; the decision procedure of Kornhauser et al.\ is a generalisation
of that parity argument to arbitrary graphs. Second, the solutions of the
theorem are long, $\bigO{\abs{V}^3}$ moves executed largely one at a time,
and nobody would fly them. Feasibility tells you that a solution exists and
that a planner which fails on a solvable instance, as prioritized planning
does in \cref{ch:ch08}, is failing for reasons of its own.

\subsection{Optimality is hard}

\begin{theorem}[Optimal MAPF is NP-hard]
\label{thm:ch07-nphard}
Given a classical MAPF instance and an integer $B$, deciding whether a
solution with $\makespan \le B$ exists is NP-hard, and so is deciding whether
a solution with $\sumcost \le B$
exists~\cite{surynek2010optimization,yu2013structure}. Unless
$\mathrm{P} = \mathrm{NP}$, no algorithm finds optimal solutions for either
objective in time polynomial in the size of the
instance.\index{MAPF!NP-hardness}\index{NP-hardness!MAPF}
\end{theorem}
\begin{proof}[Proof idea]
Both results are reductions from Boolean satisfiability.
Surynek~\cite{surynek2010optimization} proved that the optimisation variant
of multi-robot path planning on graphs is intractable; Yu and
LaValle~\cite{yu2013structure} gave reductions from 3-SAT for the makespan,
for the total arrival time, which is our sum of costs, and for the total
distance. The reductions share one idea. Every variable of the formula
becomes an agent with exactly two shortest routes, one standing for
\emph{true} and one for \emph{false}. Every clause becomes a gadget that the
routes of its literals pass through, timed so that the plan can stay at the
naive lower bound of \cref{thm:ch07-bounds} only if at least one literal of
the clause is true; when all three are false, some agent has to wait or make
a detour and pays at least one extra step. A solution with cost equal to the
lower bound therefore exists if and only if the formula is satisfiable, for
the makespan as well as for the sum of costs. The gadgets that make the
routes and the time offsets work are the substance of the two papers.
\end{proof}

NP-hardness does not say that every instance is hard. It says that a solver
which guarantees optimality must, on some instances, take exponential time,
and the solvers of \cref{ch:ch09,ch:ch10,ch:ch11} are designed around three
observations about the instances that occur in practice. Most pairs of
agents never interact, so a solver should couple agents only where their
paths conflict; Standley's independence detection~\cite{standley2010finding}
and \mstar (\cref{ch:ch11}) are built on this. The optimal cost is usually
close to the naive lower bound, so a search that starts from the naive plan
and adds one constraint at a time has little ground to cover; that is \cbs.
And a solution that is a few per cent more expensive than the optimum is
often good enough, which is what bounded-suboptimal solvers such as \ecbs
exploit. The alternative of searching the joint state space directly, with
one state per configuration of all $k$ agents, is optimal but hopeless
beyond a handful of agents: the joint space has $\abs{V}^k$ configurations,
and the brute-force searches used to verify \cref{ex:ch07-grid} explore
$14^3 = 2\,744$ of them per time step for three agents on fourteen free
cells, but would face $14^{10}$ for ten.\index{joint state space}

% ---------------------------------------------------------------------
\section{Solver families}
\label{sec:ch07-solvers}

\Cref{thm:ch07-feasible,thm:ch07-nphard} split the design space. A solver
that promises optimality must in the worst case take exponential time, a
solver that promises polynomial time must give up on optimality, and the
families in \cref{fig:ch07-solver-families} and
\cref{tab:ch07-solver-families} differ in what they promise about solution
quality and in what they pay for it~\cite{stern2019mapf,felner2017search}.
\index{solver families (MAPF)}

\begin{figure}[tb]
  \centering
  \inputfigure{ch07/solver-families}
  \caption{Five families of MAPF solvers, ordered from the strongest to the
  weakest guarantee on solution quality. The chapters in brackets treat the
  algorithms in this book; the other names are pointers to the literature.}
  \label{fig:ch07-solver-families}
\end{figure}

\begin{table}[tb]
  \centering\small
  \caption{Solver families for classical MAPF. ``Complete'' means that the
  solver returns a solution whenever one exists, given enough time; the
  guarantee refers to the objective the solver is run with.}
  \label{tab:ch07-solver-families}
  \begin{tabularx}{\textwidth}{@{}L{2.1cm}L{2.6cm}L{4.2cm}Y@{}}
  \toprule
  Family & Guarantee & Representatives & Where to use it\\
  \midrule
  optimal, search-based & optimal and complete & \cbs (\cref{ch:ch09}); \mstar (\cref{ch:ch11}); ICTS \cite{sharon2013icts}; \astar in the joint space with operator decomposition and independence detection \cite{standley2010finding} & tens of agents; when the plan is flown many times or optimality is the point\\
  bounded-suboptimal & cost $\le w \cdot$ optimum, complete & \ecbs (\cref{ch:ch10}); EECBS \cite{li2021eecbs} & hundreds of agents; the default choice for a swarm planner\\
  unbounded-suboptimal & none on cost; prioritized planning is incomplete, Push-and-Rotate is complete & prioritized planning, cooperative \astar \cite{silver2005cooperative} (\cref{ch:ch08}); Push-and-Swap, Push-and-Rotate (\cref{ch:ch11}) & real-time replanning, very many agents, sparse maps\\
  reduction-based & optimal for the encoded objective, usually the makespan; complete & SAT encodings \cite{felner2017search}; integer programming \cite{yu2016optimal} (\cref{ch:ch22}) & small dense instances; when a strong external solver is at hand\\
  learned & none & PRIMAL \cite{sartoretti2019primal} & decentralised execution with partial observation; needs a validator and a fallback\\
  \bottomrule
  \end{tabularx}
\end{table}

\paragraph{Optimal search-based solvers.}
\cbs~\cite{sharon2015cbs}\index{CBS!as solver family} searches a tree of
constraints: it plans every agent alone, finds the first conflict with the
detector of \cref{sec:ch07-detection}, and branches by forbidding one of the
two agents the conflicting vertex or edge at that time. \mstar searches the
joint space but couples only the agents that actually conflict, and
ICTS\index{ICTS} enumerates vectors of path costs in increasing order and
checks whether the paths of those costs can be combined. All three are
optimal and complete, and all three are exponential in the worst case, as
\cref{thm:ch07-nphard} demands.

\paragraph{Bounded-suboptimal solvers.}
\ecbs~\cite{barer2014ecbs} relaxes \cbs with focal search on both levels and
returns a solution of cost at most $w$ times the optimum for a chosen
$w \ge 1$, typically $w = 1.1$ to $1.5$; EECBS~\cite{li2021eecbs} adds
better lower bounds and heuristics. For a drone swarm this is the family to
start from: a few per cent of extra flight time buys an order of magnitude
in the number of agents.

\paragraph{Unbounded-suboptimal solvers.}
Prioritized planning\index{prioritized planning} \cite{silver2005cooperative}
fixes an order of the agents and plans each with space-time \astar while
treating the paths of the earlier agents as moving obstacles. It is fast and
simple, but it can fail on solvable instances and its plans can be far from
optimal; \cref{ch:ch08} shows both with counterexamples. Push-and-Swap and
Push-and-Rotate~\cite{luna2011push,dewilde2014push}\index{Push-and-Rotate}
are rule-based: they move agents to their goals one by one and resolve
blockages with local manoeuvres, and Push-and-Rotate is the complete
algorithm behind \cref{thm:ch07-feasible}.

\paragraph{Reduction-based solvers.}
Fix a makespan bound $B$, encode ``is there a solution with
$\makespan \le B$?'' as a Boolean formula or an integer program with one
variable per agent, vertex and time step, hand it to a SAT or ILP solver, and
increase $B$ until the answer is yes. The first yes is makespan-optimal by
construction. \cref{ch:ch22} shows the integer-programming
encoding~\cite{yu2016optimal}.

\paragraph{Learned solvers.}
PRIMAL~\cite{sartoretti2019primal} trains a decentralised policy by
imitation of an optimal solver and by reinforcement learning; each agent
chooses its next move from a local observation. Such policies scale and
react, but they guarantee nothing, neither completeness nor collision
freedom, so a plan they produce must pass the validator of this chapter and
a classical solver must be ready as a fallback.
